\documentclass{article}


\usepackage{arxiv}

\usepackage[utf8]{inputenc} % allow utf-8 input
\usepackage[T1]{fontenc}    % use 8-bit T1 fonts
\usepackage{hyperref}       % hyperlinks
\usepackage{url}            % simple URL typesetting
\usepackage{booktabs}       % professional-quality tables
\usepackage{amsfonts}       % blackboard math symbols
\usepackage{nicefrac}       % compact symbols for 1/2, etc.
\usepackage{microtype}      % microtypography
\usepackage{lipsum}
\usepackage{graphicx}
\graphicspath{ {./images/} }


\title{Causal Representation Learning for Brain MRI}


\author{
 Natalia Glazman}

  %% \AND
  %% Coauthor \\
  %% Affiliation \\
  %% Address \\
  %% \texttt{email} \\
  %% \And
  %% Coauthor \\
  %% Affiliation \\
  %% Address \\
  %% \texttt{email} \\
  %% \And
  %% Coauthor \\
  %% Affiliation \\
  %% Address \\
  %% \texttt{email} \\


\begin{document}
\maketitle



\section{Introduction}

Causal representation learning (CRL) is a relatively novel field seeking to recover the true data-generating and causally faithful generative factors that underlie a high-dimensional system. It has been previously shown that without any inductive biases, these factors are not uniquely identifiable.

While CRL has been widely applied to synthetic and object-centred imaging datasets to recover true generative factors, applications to real imaging data have been scarce. One of the main challenges in identifying effective CRL application to real-life images with is the lack of knowledge of both the causal graphs and the generative latent features. While previous methods have addressed identifiability of generative features via multi-view or

\section{Considerations and Background}
\label{sec:headings}

\subsection{Problem definition}
We formulate the problem using a causal latent variable model. Let $\mathbf{z}$ define a set of latent causal variables such that the observed image, $\mathbf{x}:=f(\mathbf{z}, \epsilon)$ where $\epsilon$ is exogenous noise. Given that we have access to the observed image, we wish to identify the latent generative factors up to some permutation, $h(\mathbf{z})$ as well as the causal relationship matrix that defines $f$, $\mathbf{W}$.

\subsection{Identifiability levels}
\begin{enumerate}
    \item \textbf{Component-wise identifiability}
    Component-wise identifiability of causal factors is achieved when each factor has been uniquely identified so that the learned representation $z$

    \item \textbf{Block-wise identifiability}
    A set $z$ of the latent variables is block-identified if the learned representation $z := g(x)$ contains all and only information about the ground truth $z$, i.e. $z = h(z)$ for some diffeomorphism $h : \mathbb{R} \mapsto \mathbb{R}$.
\end{enumerate}

\paragraph{Why is a block-identified representation useful for brain MRI?}

\section{Model choices}
Multi-scale, abstraction, discrete/continuous, spatial/feature vector...

\subsection{Entropy maximization}
We make the assumption here that each of the latent random variables, $z_i \in \mathbb{R}$ is one-dimensional and that the function $f$ is injective.

\subsection{Reconstruction}



\subsection{Open questions}
\begin{enumerate}
    \item Is pursuing component-wise identifiability in brain MRI worthwhile and realistic?
    \item Is block-wise identifiability worthwhile in MRI?
    \item Can we achieve good empirical results on causal factor recovery and causal graph recovery when not all theoretical guarantees are preserved?
    \item Is reconstruction necessary to preserve a faithful representation of an MRI image?
    \item Is obtaining more accurate counterfactuals possible if we can't achieve component-wise identifiability?
    \item Can we identify a causal graph without component-wise guarantees?
    \item Does having fine spatial features improve the causal estimation of the graph/block significantly?

\end{enumerate}

\section{Routes to identifiability}

Throughout this section, we invoke the invariance principle, formalised by Yao et al. to unify the different approaches to identifiability in causal representation learning.

\subsection{Multi-view}

\subsection{Soft interventions/multiple environments}

\subsection{Temporal}

\subsection{Multi-modal}

\section{Synthetic Dataset}

\subsection{Synthetic paired-contrast data}
\label{sec:synthetic-data}

Since real paired MRI has no ground-truth latents, we evaluate on a procedural
generator that renders a T1-like and a FLAIR-like $64^3$ volume of the same
subject. Each subject has a content vector $\mathbf{c}\in[-1,1]^{9}$ shared by
both views and a style vector $\mathbf{s}^{(v)}\in\mathbb{R}^{3}$ drawn
independently for each view.

\paragraph{Content.}
The nine factors are brain size, ventricle size, lesion position (3D),
cortical thickness, temporal-lobe atrophy, left--right asymmetry and sulcal
depth. They are sampled from a nonlinear structural causal model over a random
DAG, so the factors are dependent, and are then mapped monotonically to
$[-1,1]$. Each factor deforms a tissue-label map of white matter, grey matter,
ventricles and a white-matter lesion. Local factors are constructed so that
they do not change total brain volume, which keeps them distinguishable from
brain size.

\paragraph{Style.}
Each view assigns intensities to the tissue labels with a modality-specific
lookup table (e.g.\ white matter brighter than grey matter in T1, and the
reverse in FLAIR; the lesion is hypointense in T1 and hyperintense in FLAIR).
The style vector sets a global intensity gain, an offset and the noise level.
A smooth bias field and Rician noise are then applied. Since style is redrawn
for each view, it differs within every positive pair, whereas content is
identical.

\section{Evaluations}

\section{Results}
Here we show the results





\bibliographystyle{unsrt}
%\bibliography{references}  %%% Remove comment to use the external .bib file (using bibtex).
%%% and comment out the ``thebibliography'' section.





\end{document}
